\section{Introduction}

Unmanned aerial vehicle (UAV) swarms have become increasingly important in cooperative reconnaissance, target interception, area patrol, and air-combat related missions. Compared with single-platform decision making, swarm-level missions require multiple vehicles to share information, maintain safe spatial relations, and form coordinated behaviors under dynamic environmental constraints. In cooperative pursuit scenarios, several pursuer UAVs need to intercept a maneuvering target while avoiding inter-UAV collisions and mission timeouts. When communication is limited or the coordination policy is unstable, the team may suffer from inefficient pursuit, unsafe convergence, or failed interception.

Multi-agent reinforcement learning (MARL) provides an end-to-end learning framework for cooperative UAV decision making. The centralized training and decentralized execution paradigm allows a critic to exploit global information during training while each agent executes its policy based on local observations. MAPPO has been shown to be a strong and practical baseline for cooperative multi-agent tasks~\cite{yu2021mappo}. However, a standard MAPPO policy lacks an explicit mechanism for relational reasoning. Under a limited communication radius, relative spatial relations and communication reachability among UAVs may not be fully exploited by a policy that mainly relies on local observations and a centralized critic.

Graph neural networks and graph attention mechanisms provide a natural way to model relations among multiple agents~\cite{velickovic2017gat,malysheva2020magnet,liu2024gnn_marl}. Nevertheless, standard graph attention usually computes attention scores from node embeddings. In UAV pursuit tasks, edge-level physical relations such as relative distance, bearing, relative velocity, and communication reachability directly affect formation geometry, collision avoidance, and target encirclement. If these relations are only implicitly encoded in node states, the policy must learn to infer them indirectly, which can increase training difficulty and reduce robustness when the communication topology changes.

Realistic UAV swarms cannot assume global and stable communication. Communication radius, link quality, bandwidth, and occlusion may all change the set of information available to each platform. Existing learning-to-communicate and graph-based MARL studies have investigated when and how agents should exchange information~\cite{singh2018ic3net,cuzin2026gnn_comm_survey}, and UAV communication or mobile edge computing studies have also combined graph attention with multi-agent learning~\cite{feng2024gat_uav_comm,kim2024uav_mec_madrl}. In cooperative pursuit, however, communication constraints are not only graph connectivity changes. The geometry and motion relationship between agents also determine whether an information edge is useful for safe and efficient pursuit.

This paper focuses on a realistic and implementable research step: validating a limited-communication graph MARL method in a simplified two-dimensional heterogeneous UAV pursuit environment before moving to full 6DOF air combat, missile, radar, and human-UAV teaming systems. This choice reduces simulation complexity, isolates the algorithmic contribution, and provides a reusable representation layer for future migration to LAG/JSBSim or other 6DOF simulators. We propose EA-RG-MAPPO-S, an edge-aware role graph MAPPO method with staged random-radius fine-tuning.

The main contributions are as follows:
\begin{enumerate}
    \item We propose an edge-aware role graph policy representation for limited-communication UAV pursuit. Pursuers and targets are modeled as dynamic graph nodes with role semantics, and relative position, distance, bearing, relative velocity, and communication reachability are injected into the graph attention score.
    \item We introduce a staged random-radius fine-tuning strategy. The policy first learns basic coordination under a fixed communication radius and is then fine-tuned with randomly sampled radii, improving stability across communication topologies instead of overfitting to a single radius.
    \item We build a reproducible heterogeneous UAV pursuit benchmark and evaluate MAPPO, GAT-MAPPO, EA-RG-MAPPO-S, and ablation variants over multiple seeds. The final 300-episode-per-seed evaluation shows that EA-RG-MAPPO-S keeps collision rates between 0.054 and 0.086 across four communication radii.
\end{enumerate}

It should be noted that an auxiliary target-intent prediction branch was explored during development. Diagnostic results show insufficient balanced accuracy, indicating that this branch cannot yet support a high-accuracy intent-recognition claim. Therefore, the main contribution of this paper is limited-communication edge-aware role graph coordination rather than target intent recognition.
